\begin{apartats}

\begin{tria}{domini-funcions}
\itemtria{original}{0,75}{1,25}
Determina el domini i estudia la continuïtat de les funcions següents:
\begin{graella}{2}
  \sa y=\sqrt{x^2-9} & \sa y=\ln(4-x)
\end{graella}

\begin{solucio}
i) Cal $x^2-9\ge0$, és a dir $|x|\ge3$: $\mathrm{Dom}=(-\infty,-3]\cup[3,+\infty)$.
És contínua a tot el domini (composició de contínues).\\
ii) Cal $4-x>0$: $\mathrm{Dom}=(-\infty,4)$. És contínua a tot el domini.
\end{solucio}

\itemtria{signes-i-condicions}{0,75}{1,25}
Determina el domini i estudia la continuïtat de les funcions següents:
\begin{graella}{2}
  \sa y=\sqrt{\frac{x-1}{x+2}} & \sa y=\frac{x+1}{\ln x}
\end{graella}

\begin{solucio}
i) Cal $\dfrac{x-1}{x+2}\ge0$, amb $x\ne-2$. Taula de signes: el quocient és positiu a
$(-\infty,-2)$, negatiu a $(-2,1)$, nul en $x=1$ i positiu a $(1,+\infty)$:
$\mathrm{Dom}=(-\infty,-2)\cup[1,+\infty)$. És contínua a tot el domini.\\
ii) Cal $x>0$ (logaritme) i $\ln x\ne0$, és a dir $x\ne1$: $\mathrm{Dom}=(0,1)\cup(1,+\infty)$.
És contínua a tot el domini. En $x=1$ els laterals valen $-\infty$ i $+\infty$:
\textbf{salt infinit}.
\end{solucio}
\end{tria}

\begin{tria}{discontinuitats-racional}
\itemtria{original}{1}{1,25}
Troba els punts en què la funció
\[
f(x)=\frac{x^2-x-6}{x^2-2x-3}
\]
és discontínua i classifica'n la discontinuïtat.

\begin{solucio}
$f(x)=\dfrac{(x-3)(x+2)}{(x-3)(x+1)}=\dfrac{x+2}{x+1}$ per a $x\ne3$.
El domini és $\mathbb{R}\setminus\{-1,3\}$.\\
$x=3$: $\lim_{x\to3}f(x)=\dfrac54$ existeix però $f(3)$ no. \textbf{Discontinuïtat evitable.}\\
$x=-1$: els laterals valen $-\infty$ i $+\infty$. \textbf{Salt infinit} (asímptota vertical $x=-1$).
\end{solucio}

\itemtria{factor-doble}{1}{1,25}
Troba els punts en què la funció
\[
h(x)=\frac{x^2-x-6}{x^2-6x+9}
\]
és discontínua i classifica'n la discontinuïtat.

\begin{solucio}
$h(x)=\dfrac{(x-3)(x+2)}{(x-3)^2}$: el domini és $\mathbb{R}\setminus\{3\}$.\\
El factor $(x-3)$ és comú, però en simplificar encara en queda un al denominador:
$h(x)=\dfrac{x+2}{x-3}$ per a $x\ne3$. En $x=3$ el numerador tendeix a $5$ i el denominador a
$0$: els laterals valen $-\infty$ (per l'esquerra) i $+\infty$ (per la dreta).
\textbf{Salt infinit}, no evitable: un factor comú no fa evitable la discontinuïtat si el
denominador encara s'hi anul·la després de simplificar.
\end{solucio}
\end{tria}

\begin{nomesllarg}
\apartat{0,75}
\textbf{Inventa.} Escriu una funció racional $g$ tal que
\[
\lim_{x\to+\infty}g(x)=3
\]
i que tingui una discontinuïtat de salt infinit en $x=2$ i una discontinuïtat evitable en
$x=-1$.

\begin{solucio}
Per exemple $g(x)=\dfrac{3(x+1)(x-5)}{(x+1)(x-2)}$.\\
El factor $(x+1)$ es cancel·la: discontinuïtat evitable en $x=-1$, amb límit
$\frac{3(-6)}{-3}=6$. En $x=2$ el denominador s'anul·la i el numerador no: salt infinit.
Numerador i denominador tenen el mateix grau i el quocient dels coeficients principals
és $3$, així que $\lim_{x\to+\infty}g(x)=3$. (La resposta no és única.)
\end{solucio}
\end{nomesllarg}

\end{apartats}
